\documentclass[11pt]{article}
\usepackage[margin=1in]{geometry}
\usepackage{amsmath}
\usepackage{amssymb}
\usepackage{hyperref}
\usepackage[numbers]{natbib}
\usepackage{booktabs}
\usepackage{multirow}
\hypersetup{colorlinks=true, linkcolor=blue, citecolor=blue, urlcolor=blue}

\title{Daily Paper Review: TESSY --- Teacher-Student Cooperation\\for Reasoning Model Fine-Tuning}
\author{Internbot --- Automated Research Review}
\date{April 20, 2026}

\begin{document}
\maketitle

\section{Paper Summary}

\textbf{Title:} How to Fine-Tune a Reasoning Model? A Teacher-Student Cooperation Framework to Synthesize Student-Consistent SFT Data \\
\textbf{Authors:} Zixian Huang, Kaichen Yang, Xu Huang, Feiyang Hao, Qiming Ge, Bowen Li, He Du, Kai Chen, Qipeng Guo (Shanghai AI Lab) \\
\textbf{arXiv:} \href{https://arxiv.org/abs/2604.14164}{2604.14164} \\
\textbf{Topics:} LLM Efficiency, Knowledge Distillation, Reasoning Models

\subsection{Problem}

A widely adopted strategy for improving LLMs is supervised fine-tuning (SFT) on data synthesised by a stronger teacher model. For \textbf{reasoning models} (e.g., Qwen3-8B), this standard approach not only fails but actively \textit{degrades} performance. Using GPT-OSS-120B as the teacher: LiveCodeBench-Pro drops by $-3.25\%$ and OJBench drops by $-10.02\%$. This is counterintuitive---the teacher is stronger and its data is higher quality, yet training on it hurts the student.

TESSY identifies \textbf{stylistic divergence} as the root cause. Unlike base or instruct models, reasoning models have undergone extensive post-training (RL, large-scale SFT) that has entrenched distinctive stylistic patterns---discourse markers (``Okay, let's see,'' ``But wait''), reasoning flow structure, self-correction patterns, and word frequency profiles. The SFT loss decomposes into:
\begin{equation}
    \mathcal{L}(\mathcal{M}_S) = \underbrace{\sum_{i \in \mathcal{T}_{\mathrm{Cap}}} \mathbb{E}\!\left[\mathrm{KL}\!\left(P_D(y_i) \| P_{\mathcal{M}_S}(y_i)\right)\right]}_{\mathcal{L}_{\mathrm{Cap}}} + \underbrace{\sum_{i \in \mathcal{T}_{\mathrm{Sty}}} \mathbb{E}\!\left[\mathrm{KL}\!\left(P_D(y_i) \| P_{\mathcal{M}_S}(y_i)\right)\right]}_{\mathcal{L}_{\mathrm{Sty}}}
\end{equation}
For base models, $\mathcal{L}_{\mathrm{Sty}}$ is harmless because the model has no strong stylistic priors. For reasoning models, learning the teacher's alien style creates interference that degrades the capability loss $\mathcal{L}_{\mathrm{Cap}}$. The authors verify this with a gradient of interventions: the more the teacher participates in data synthesis, the worse the degradation (Teacher-Only $-10.02\%$, Teacher-Mix $-6.90\%$, Teacher-Think $-6.25\%$, Teacher-Reference $-4.53\%$, Teacher-Score $-0.43\%$ on OJBench).

\subsection{Method}

TESSY (Teacher-Student Cooperation Data Synthesis) disentangles capability from style by interleaving teacher and student generation at the token span level.

\paragraph{Alternating Generation.}
TESSY generates each training response as an alternating sequence $y = [s^1, t^1, s^2, t^2, \ldots]$, where $s^i$ are \textbf{student spans} (style tokens) and $t^i$ are \textbf{teacher spans} (capability tokens). Each span is conditioned on all previously generated spans:
\begin{align}
    s^i &= \mathcal{M}_S\!\left(x,\, [s^1, t^1, \ldots, s^{i-1}, t^{i-1}]\right) \\
    t^i &= \mathcal{M}_T\!\left(x,\, [s^1, t^1, \ldots, s^i]\right)
\end{align}
Generation starts with the student because model outputs typically begin with stylistic phrases.

\paragraph{Generate-Then-Rollback.}
The key technical challenge is deciding when to switch between teacher and student. TESSY uses a generate-then-rollback strategy: (1) the current model generates $k = 20$ tokens, (2) a \textbf{boundary predictor} (token-level binary classifier on Qwen3-0.6B-Base) classifies each token as style or capability, identifying the first token that violates the expected type, (3) all tokens from the boundary onward are discarded, and (4) generation switches to the other model. Formally:
\begin{equation}
    t^i = \tilde{t}^i_{<B_T(\tilde{t}^i)}, \qquad s^i = \tilde{s}^i_{<B_S(\tilde{s}^i)}
\end{equation}
where $B_T$ and $B_S$ are the capability and style boundary predictors, respectively.

\paragraph{Style vs.\ Capability Tokens.}
\textbf{Style tokens} include discourse markers (``well,'' ``okay,'' ``let's see''), transitions, hesitation, and narration that does not contain reasoning. \textbf{Capability tokens} include code, formulas, numbers, and actual deductive reasoning steps. The boundary predictors are trained on 100k segments annotated by the teacher model (GPT-OSS-120B) using a structured prompt.

\paragraph{Final Answer Delegation.}
The entire final answer (outside \texttt{<think>} tags) is always generated by the \textbf{student}. Using the teacher for final answers caused drops of $-12.33\%$ and $-13.58\%$, confirming that ``avoiding style conflicts is more important than achieving marginal improvements in data quality.''

\paragraph{Implementation.}
Built on vLLM with prefix caching. Approximately 77.65\% of tokens in TESSY output come from the teacher, 22.4\% from the student. Despite majority-teacher composition, the output distribution is student-aligned (TF-IDF similarity to student increases from 0.849 to 0.898 for GPT-OSS, and from 0.846 to 0.960 for DS-R1). TESSY responses are also shorter than teacher-only responses by 7,594--8,938 tokens on average, because the student's earlier termination tendency guides the teacher.

\subsection{Results}

\begin{table}[h]
\centering
\caption{Code generation results. Student: Qwen3-8B, Teacher: GPT-OSS-120B. 9-epoch full fine-tuning.}
\begin{tabular}{lcccc}
\toprule
\textbf{Method} & \textbf{LCB-V5} & \textbf{LCB-V6} & \textbf{LCB-Pro} & \textbf{OJBench} \\
\midrule
Qwen3-8B (baseline) & 55.09 & 49.58 & 25.35 & 18.75 \\
+ Teacher-Only & 41.32 & 36.57 & 22.10 & 8.73 \\
+ Teacher-Score & 54.55 & 48.00 & 26.06 & 18.32 \\
\textbf{+ TESSY} & \textbf{62.87} & \textbf{55.43} & \textbf{36.69} & \textbf{25.43} \\
\bottomrule
\end{tabular}
\end{table}

Key findings:

\begin{itemize}
    \item \textbf{TESSY turns degradation into improvement:} On OJBench, Teacher-Only causes $-10.02\%$; TESSY achieves $+6.68\%$. On LCB-Pro, Teacher-Only causes $-3.25\%$; TESSY achieves $+11.34\%$. A swing of 16.70 and 14.59 percentage points, respectively.

    \item \textbf{Out-of-domain generalisation:} TESSY trained on code data generalises to math: $+3.75\%$ on AIME2024 and $+0.93\%$ on AIME2025, while Teacher-Only degrades AIME2025 by $-2.71\%$ and GPQA by $-4.35\%$.

    \item \textbf{Cross-family advantage scales with distance:} The more stylistically distant the teacher, the larger TESSY's advantage. TESSY vs.\ Teacher-Only: $+1.07\%$ for same-family Qwen3-235B, $+3.45\%$ for DeepSeek-R1, $+16.79\%$ for cross-family GPT-OSS-120B.

    \item \textbf{Cross-training fails:} Training Qwen3-8B on TESSY data synthesised with Qwen3-30B-A3B as the student causes $-1.94\%$ on OJBench. Even within the same model family, \textit{the specific student's style matters}.

    \item \textbf{TESSY output quality exceeds teacher:} Despite 22.4\% student tokens, TESSY output evaluated directly (not as training data) outperforms GPT-OSS-120B by $+10.99\%$ on OJBench at matched generation budget. The student's presence improves efficiency.

    \item \textbf{LoRA insufficient:} LoRA (rank 8, $\alpha = 16$) performed significantly worse than full fine-tuning, even worse than Teacher-Only with full parameters. Reasoning models require substantial parameter updates.
\end{itemize}

\subsection{Limitations}

\begin{itemize}
    \item \textbf{Capability ceiling:} When generation length is unrestricted, Teacher-Only slightly outperforms TESSY, suggesting the student's influence may cause premature reasoning termination.

    \item \textbf{Boundary predictor accuracy:} No explicit accuracy metrics are reported for the style/capability classifiers. The framework's quality depends on this classification.

    \item \textbf{Code-centric evaluation:} Primary experiments are on code generation. Out-of-domain math results are encouraging but broader task coverage (creative writing, multi-turn dialogue) is unexplored.

    \item \textbf{No combination with reject sampling:} TESSY isolates the distribution alignment effect but does not combine with data quality techniques that could provide further gains.

    \item \textbf{Fixed $k = 20$:} The span length is set empirically with no sensitivity analysis. Optimal values may vary by task or model pair.
\end{itemize}

\subsection{Relevance to Our Research}

TESSY provides the missing ``distribution alignment'' piece for distillation in reasoning models, directly complementing several threads in our review series.

\textbf{Resolving the self-distillation paradox:} Our review of ``Why Does Self-Distillation (Sometimes) Degrade Reasoning?''~\cite{selfdistill} identified that self-distillation can hurt because the student learns to imitate superficial patterns rather than reasoning structure. TESSY provides a mechanistic explanation: the damage comes specifically from \textit{style tokens}, not capability tokens. Self-distillation avoids cross-model style divergence (same model $\Rightarrow$ same style) but introduces a different problem: the teacher's errors are stylistically indistinguishable from correct reasoning. TESSY's capability/style decomposition suggests a refined self-distillation where capability tokens are verified against ground truth while style tokens are preserved unconditionally.

\textbf{Connection to xLSTM Distillation}~\cite{xlstm}: xLSTM distillation transfers knowledge across architectures (Transformer $\to$ xLSTM), facing both architectural and stylistic mismatch. TESSY's framework could be adapted: use the Transformer teacher for capability tokens (attention patterns, reasoning steps) while letting the xLSTM student generate style tokens that match its recurrent architecture's natural output distribution.

\textbf{Connection to RLSD}~\cite{rlsd}: RLSD converts binary RL rewards into dense token-level credit via evidence reweighting. TESSY's style/capability decomposition provides a complementary signal: style tokens should receive zero credit in RL (they do not affect task correctness), while capability tokens should receive full credit. Combining TESSY's token classification with RLSD's evidence reweighting could yield a more precise credit assignment that ignores stylistic variance and focuses gradient signal on reasoning-relevant tokens.

\textbf{Connection to KnowRL}~\cite{knowrl}: Both papers address the question of \textit{what to inject} during training. KnowRL injects minimal-sufficient knowledge points as prompt hints; TESSY injects teacher capability tokens directly into the generation stream. KnowRL operates at the problem level (which hints to include), TESSY at the token level (which tokens to take from the teacher). A unified framework could use KnowRL's CSS to select which problems need teacher augmentation, then TESSY to perform the augmentation at token granularity.

\section{Follow-up Research Directions}

\subsection{Direction 1: TESSY-Guided Credit Assignment for RL Post-Training}

\textbf{Gap:} RLSD~\cite{rlsd} provides dense token-level credit via evidence reweighting, but treats all tokens uniformly---style tokens receive credit even though they do not affect task correctness. KnowRL~\cite{knowrl} showed that \textit{what} the model sees matters; TESSY shows that \textit{which tokens} carry capability matters. No existing RL method uses style/capability decomposition for credit assignment.

\textbf{Approach:}
\begin{enumerate}
    \item Train TESSY's boundary predictor on the RL policy's own outputs (not the teacher's). This produces a token-level binary mask $m_i \in \{0, 1\}$ for each position, where $m_i = 1$ indicates a capability token.

    \item Modify RLSD's evidence reweighting to incorporate the mask. The per-token advantage becomes:
    \begin{equation}
        \hat{A}_i = m_i \cdot w_i \cdot (R - b)
    \end{equation}
    where $w_i$ is RLSD's evidence weight and $R - b$ is the baseline-subtracted reward. Style tokens ($m_i = 0$) receive zero gradient, focusing all learning signal on reasoning-relevant tokens.

    \item Apply this within GRPO training on math benchmarks. Compare against vanilla GRPO, RLSD, and KnowRL. The key hypothesis: masking style tokens reduces gradient noise from stylistic variance across rollouts, improving sample efficiency.

    \item Extend to RAGEN-2's~\cite{ragen2} SNR-aware filtering: compute SNR separately for capability tokens and style tokens. Problems where capability-token SNR is high but overall SNR is low (due to style variance) should be retained, not filtered.
\end{enumerate}

\textbf{Feasibility:} High. The boundary predictor is lightweight (Qwen3-0.6B-Base, already trained). Modifying RLSD's per-token weighting requires only a multiplicative mask---no architectural changes. Start with Qwen3-8B on GSM8K/MATH using GRPO, comparing masked vs.\ unmasked credit. The target: matching RLSD's improvement with 30--50\% fewer training steps by eliminating gradient noise from style tokens.

\subsection{Direction 2: Style-Preserving Distillation for Diffusion LLMs}

\textbf{Gap:} LangFlow~\cite{langflow} showed that continuous diffusion can rival discrete for language modeling, and D-MMD~\cite{dmmd} distils discrete DLMs via MMD matching. Neither addresses the style divergence problem: when distilling a large diffusion teacher into a smaller student, the student may inherit the teacher's denoising style (noise schedule preferences, token generation order) rather than its reasoning capability. TESSY's insight---that style and capability are separable---has not been explored for diffusion model distillation.

\textbf{Approach:}
\begin{enumerate}
    \item Define ``style'' and ``capability'' for diffusion LLMs. In the denoising process, \textbf{capability} corresponds to which tokens are correctly predicted at each step (accuracy of the denoiser), while \textbf{style} corresponds to the order in which tokens are revealed and the confidence calibration across positions.

    \item For D-MMD distillation, apply TESSY's boundary predictor to the \textit{final generated text} and compute the MMD loss only on capability-token positions. Style-token positions use a student-consistency loss (KL between student's own predictions at adjacent steps) instead of teacher matching.

    \item For LangFlow's ODE-based distillation via Consistency Models, apply the style/capability mask to the trajectory: match the teacher's ODE trajectory at capability-token embedding dimensions while allowing the student to maintain its own trajectory at style-token dimensions.

    \item Compare against standard D-MMD and Consistency Model distillation at matched step budgets. The hypothesis: style-preserving distillation yields better downstream task performance because the student's denoising patterns remain internally consistent.
\end{enumerate}

\textbf{Feasibility:} Medium. The main challenge is defining style vs.\ capability in the diffusion denoising space, which is less clear-cut than in autoregressive generation. Start with D-MMD on MDLM-130M (where the text-level boundary predictor is directly applicable to final outputs), then extend to LangFlow's continuous setting. Use LM1B for evaluation, targeting Gen.PPL improvement over standard D-MMD at 4--8 denoising steps.

\subsection{Direction 3: Adaptive Style--Capability Boundaries via Curriculum Learning}

\textbf{Gap:} TESSY uses a fixed boundary predictor trained before SFT begins. As the student learns from teacher capability tokens, its own capability improves and what counts as ``style'' vs.\ ``capability'' may shift. Early in training, the student may classify complex reasoning steps as capability (needing teacher help); late in training, those same steps may be within the student's capability and should be classified as style (to preserve the student's own improved reasoning flow). KnowRL~\cite{knowrl} faced an analogous problem: static KP selection ignores policy improvement during training.

\textbf{Approach:}
\begin{enumerate}
    \item Re-train the boundary predictor every $T_{\mathrm{eval}}$ epochs using the \textit{current} student checkpoint's outputs. As the student improves, more tokens shift from ``capability'' (needing teacher) to ``style'' (student-sufficient), progressively increasing the student's contribution.

    \item Define a curriculum schedule: at epoch $e$, the teacher contribution ratio is $\rho(e) = \rho_0 \cdot (1 - e / E_{\mathrm{max}})^\beta$, where $\rho_0 \approx 0.78$ (TESSY's initial 77.65\%) and $\beta$ controls the decay rate. This creates a withdrawal schedule analogous to Direction~1 in our KnowRL review (progressive knowledge withdrawal).

    \item Monitor the cross-training diagnostic: periodically evaluate whether the current TESSY data would help a \textit{different} student (e.g., Qwen3-30B-A3B). If the data transfers well, the teacher contribution may be too dominant (style leakage); if it transfers poorly, the student contribution is appropriate. This self-diagnostic prevents excessive teacher dependence.

    \item Combine with TESSY's observation that LoRA is insufficient: use full fine-tuning for early epochs (when the distribution shift is largest) and switch to LoRA for late epochs (when the student's distribution has stabilised and only incremental capability gains are needed).
\end{enumerate}

\textbf{Feasibility:} High. Re-training a 0.6B boundary predictor is cheap ($< 1\%$ of SFT compute). The curriculum schedule adds no overhead beyond periodic evaluation. Start with Qwen3-8B + GPT-OSS-120B on the existing TESSY-Code-80K dataset, comparing fixed vs.\ adaptive boundaries over 9 epochs. The target: reducing TESSY's capability ceiling (the gap where Teacher-Only wins at unlimited generation length) while preserving the distribution alignment advantage.

\section{Conclusion}

TESSY resolves a fundamental tension in reasoning model distillation: teacher data is higher quality but distributionally incompatible with the student. By decomposing generation into style tokens (from the student) and capability tokens (from the teacher) via a generate-then-rollback framework, TESSY converts a $-10.02\%$ degradation into a $+6.68\%$ improvement on OJBench---a 16.70-point swing. The most striking finding is that 77.65\% teacher tokens can coexist with student-aligned distribution, and that the student's presence actually \textit{improves} generation efficiency by guiding the teacher toward earlier termination.

For our lab, TESSY fills a critical gap in the distillation pipeline. Our previous reviews of self-distillation~\cite{selfdistill} and xLSTM distillation~\cite{xlstm} identified that distillation can fail, but lacked a mechanistic explanation. TESSY provides one: the damage comes from style tokens, not capability tokens. This insight is immediately actionable: any distillation or RL method that processes token-level signals (RLSD's evidence reweighting, KnowRL's knowledge injection, RAGEN-2's SNR filtering) can benefit from masking out style tokens to reduce gradient noise.

The broader principle---that reasoning models have entrenched stylistic priors that resist overwriting---has implications beyond distillation. It suggests that RL post-training (GRPO, PPO) may also suffer from style interference: reward signals from rollouts that happen to match the model's stylistic preferences may be confounded with genuine capability improvements. Direction~1 proposes using TESSY's boundary predictor to separate these signals, potentially improving the sample efficiency of all RL methods in our review series.

Across 34 reviews, a meta-pattern emerges: the most effective training methods respect the model's existing structure rather than overwriting it. KnowRL injects minimal knowledge to avoid dependency; TRACE~\cite{trace} routes to per-capability adapters to avoid interference; and now TESSY preserves the student's style to avoid distribution conflict. The principle is consistent: \textbf{targeted augmentation outperforms wholesale replacement}.

\begin{thebibliography}{10}
\bibitem{tessy} Huang, Z., et al. (2026). How to Fine-Tune a Reasoning Model? A Teacher-Student Cooperation Framework to Synthesize Student-Consistent SFT Data. \textit{arXiv:2604.14164}.
\bibitem{selfdistill} Jiao, W., et al. (2026). Why Does Self-Distillation (Sometimes) Degrade the Reasoning Capability of LLMs? \textit{arXiv:2603.24472}.
\bibitem{xlstm} Beck, M., et al. (2026). Effective Distillation to Hybrid xLSTM Architectures. \textit{arXiv:2603.15590}.
\bibitem{rlsd} Yang, C., et al. (2026). Self-Distilled RLVR. \textit{arXiv:2604.03128}.
\bibitem{knowrl} Yu, L., et al. (2026). KnowRL: Boosting LLM Reasoning via RL with Minimal-Sufficient Knowledge Guidance. \textit{arXiv:2604.12627}.
\bibitem{ragen2} Wang, H., et al. (2026). RAGEN-2: Reasoning Collapse in Agentic RL. \textit{arXiv:2604.06268}.
\bibitem{trace} Kang, H., et al. (2026). TRACE: Capability-Targeted Agentic Training. \textit{arXiv:2604.05336}.
\bibitem{langflow} Chen, Y., et al. (2026). LangFlow: Continuous Diffusion Rivals Discrete in Language Modeling. \textit{arXiv:2604.11748}.
\bibitem{dmmd} Schiff, Y., et al. (2026). Beyond Single Tokens: Distilling Discrete Diffusion Models via Discrete MMD. \textit{arXiv:2603.20155}.
\bibitem{oaem} Kennedy, I. W. \& Moosavi, N. S. (2026). Initialisation Determines the Basin: Efficient Codebook Optimisation for Extreme LLM Quantization. \textit{arXiv:2604.08118}.
\end{thebibliography}

\end{document}
